
% 参考文献
% 优先引用中文文献,8-15 条,GB/T 7714-2015 格式,每条在正文有引用。
%
% 对于常见的各类参考文献标注方法如下:
%   1) 著作:作者姓名,题名[M].出版地:出版社,出版年.
%   2) 期刊论文:作者姓名.题名[J].期刊名称,年,卷(期):页码.
%   3) 会议论文集:作者姓名.题名[C]//论文集名称,会议地点,会议日期.
%   4) 学位论文:作者姓名.题名[D].出版地:出版者,出版年.
%   5) 专利文献:专利申请者或所有者姓名.专利题名:专利国别,专利号[P].公告日期或公开日期.获取路径.
%   6) 电子文献:作者姓名.题名[文献类型标志(含文献载体标志)见其它].出版地:出版者,出版年(更新或修改日期),获取路径.
%   7) 报告:作者姓名.题名[R].出版地:出版者,出版年.
%   8) 标准:标准号.题名[S].出版地:出版者,出版年.
%
% 注意使用 AI 工具生成的内容,应在正文相应位置进行标注;
% 在参考文献中列出所用 AI 工具,格式如下:
%   [编号]工具名称,版本/型号,开发机构/公司,使用日期。
%   示例:[1]DeepSeek,DeepSeek-R1-0528,深度求索(DeepSeek),2025-09-05

\begin{thebibliography}{99}

% ========== 中文文献(优先) ==========
\bibitem{ref1} 【作者1】, 【作者2】. 【文献标题】[J]. 【期刊名】, 【年】, 【卷(期)】: 【起止页】.

\bibitem{ref2} 【作者1】, 【作者2】. 【文献标题】[M]. 【出版社所在地】: 【出版社】, 【年】.

\bibitem{ref3} 【作者1】, 【作者2】. 【文献标题】[D]. 【所在地】: 【授予单位】, 【年】.

\bibitem{ref4} 【作者1】, 【作者2】. 【文献标题】[C]//【会议名】. 【会议地点】: 【出版社】, 【年】: 【起止页】.

\bibitem{ref5} 【作者1】, 【作者2】. 【文献标题】[R]. 【所在地】: 【机构】, 【年】.

\bibitem{ref6} 【作者1】, 【作者2】. 【文献标题】[J]. 【期刊名】, 【年】, 【卷(期)】: 【起止页】.

\bibitem{ref7} 【作者1】, 【作者2】. 【文献标题】[J]. 【期刊名】, 【年】, 【卷(期)】: 【起止页】.

\bibitem{ref8} 【作者1】, 【作者2】. 【文献标题】[M]. 【出版社所在地】: 【出版社】, 【年】.

\bibitem{ref9} 【作者1】, 【作者2】. 【文献标题】[D]. 【所在地】: 【授予单位】, 【年】.

\bibitem{ref10} 【作者1】, 【作者2】. 【文献标题】[J]. 【期刊名】, 【年】, 【卷(期)】: 【起止页】.

% ========== AI 工具(必列) ==========
% 格式:[编号]工具名称,版本/型号,开发机构/公司,使用日期。
\bibitem{aideepseek} DeepSeek, DeepSeek-R1-0528, 深度求索(DeepSeek), 2025-09-05.

\end{thebibliography}
